(1) 平面界面切向波矢守恒给 $n_1\sin\theta_i=n_2\sin\theta_t$。在电场振幅的常用 p 偏振约定下，Fresnel 反射系数可写 $r_p=(n_2\cos\theta_i-n_1\cos\theta_t)/(n_2\cos\theta_i+n_1\cos\theta_t)$；不同选择 p 偏振基向可能改变整体符号，但零点条件不变。
(2) p 偏振零反射要求 $n_2\cos\theta_i=n_1\cos\theta_t$。与 $n_1\sin\theta_i=n_2\sin\theta_t$ 联立，得到 $\theta_i+\theta_t=\pi/2$ 及 $\tan\theta_B=n_2/n_1$，所以 $\theta_B=\arctan(n_2/n_1)$。必须同时使用边界条件和几何折射关系。
(3) 因为反射角等于入射角，反射射线与折射射线在界面法平面内夹角为 $\pi-(\theta_B+\theta_t)=\pi/2$，所以两传播方向垂直。s 偏振系数为 $r_s=(n_1\cos\theta_i-n_2\cos\theta_t)/(n_1\cos\theta_i+n_2\cos\theta_t)$，在 $\theta_B$ 一般不为零，因此反射光偏向 s 偏振。
(4) 未偏振光的两正交偏振入射强度各半，因 $R_p(\theta_B)=0$，总反射率 $R_{\rm unpol}(\theta_B)=[|r_s(\theta_B)|^2+|r_p(\theta_B)|^2]/2=|r_s(\theta_B)|^2/2$。若介质有吸收，折射率变复数，实入射角下通常不能同时抵消 p 反射振幅的实、虚部分，因此理想零反射的布儒斯特角可能被一个非零反射最小值替代。
